\documentclass[12pt]{article}

\input{../../_assets/preambulo.tex}


\begin{document}

    % 1. Foto de fondo
    % 2. Título
    % 3. Encabezado Izquierdo
    % 4. Color de fondo
    % 5. Coord x del titulo
    % 6. Coord y del titulo
    % 7. Fecha

    
    \input{../../_assets/portada}
    \portadaExamen{ffccA4.jpg}{EDIP\\Examen III}{EDIP. Examen III}{MidnightBlue}{-8}{28}{2025-26}{David Sánchez Muñoz}

    \begin{description}
         \item[Asignatura] Estadística Descriptiva e Introducción a la Probabilidad.
        \item[Curso Académico] 2024-25.
        \item[Grado] Doble Grado en Ingeniería Informática y Matemáticas.
        \item[Grupo] Único.
        \item[Profesor] Juan Antonio Maldonado Jurado.
        \item[Descripción] Convocatoria Ordinaria.
        \item[Fecha] 06/06/2025.
        \item[Duración] 2 horas.
    \end{description}

    \newpage


    % ------------------------------------
    
    \begin{ejercicio}
        La tabla informa sobre el resultado de 20 mediciones:
        \begin{center}
        \renewcommand{\arraystretch}{1.2}
        \begin{tabular}{c|ccc}
            $X \setminus Y$ & $[0-8]$ & $\left]8-20\right]$ & $\left]20-50\right]$ \\ \hline
            $-1$ & $0$ & $0$ & $4$ \\
            $0$  & $0$ & $2$ & $3$ \\
            $1$  & $4$ & $2$ & $0$ \\
            $2$  & $5$ & $0$ & $0$
        \end{tabular}
        \end{center}

        \begin{enumerate}[label=\alph*)]
            \item Supuesto $|X| \leq 1$, determinar el número de mediciones de $Y$ cuyos valores oscilan entre el $25\%$ de los índices más altos y el valor modal.
            \item ¿Qué índice medio es más representativo: $X$ o $Y$?
            \item Estudiar la independencia lineal de las variables.
            \item Estimar mediante una recta de regresión mínimo cuadrática el valor de $Y$ cuando $X=0$. ¿Coincide este valor con la estimación de menor error cuadrático medio de $Y$ cuando $X=0$? Razone.
        \end{enumerate}
    \end{ejercicio}

    \begin{ejercicio}
        Para una determinada población se ha definido una variable aleatoria $X$ que representa el número de casos de enfermedad rara (por cada $10^3$ habitantes) registrados anualmente, con función de distribución dada por:
        \[
        F(x) = \begin{cases}
            0 & \text{si } x < 0 \\
            0.05 & \text{si } 0 \leq x < 1 \\
            0.15 & \text{si } 1 \leq x < 2 \\
            0.43 & \text{si } 2 \leq x < 3 \\
            0.65 & \text{si } 3 \leq x < 4 \\
            0.8 & \text{si } 4 \leq x < 5 \\
            k & \text{si } 5 \leq x < 6 \\
            1 & \text{si } x \geq 6
        \end{cases}
        \]

        \begin{enumerate}[label=\alph*)]
            \item Valor de $k$ sabiendo que $E[X] = 3$.
            \item Cuartiles de la distribución y varianza de $X$.
            \item Probabilidad de que un determinado año el número de casos sea por lo menos 2 si se sabe que no llega a 6.
            \item La experiencia indica que cuando el número de casos es inferior a 3, la probabilidad de tratar correctamente la enfermedad es del $5\%$, del $15\%$ cuando el número de casos está entre 3 y 4 y del $80\%$ cuando es superior a 4. Calcular la probabilidad de que en dicha población los enfermos hayan recibido el tratamiento adecuado.
        \end{enumerate}
    \end{ejercicio}

    \begin{ejercicio}
        Se tiene dos dados cuyas seis caras están marcadas con los seis divisores del número 2012.

        \begin{enumerate}[label=\alph*)]
            \item ¿Cuál es la probabilidad de que al lanzar los dados la suma de las dos caras superiores también sea divisor de 2012?
            \item Si se lanzan los dados 10 veces, calcule la probabilidad de que al menos una vez, la suma de las caras superiores haya resultado ser un divisor.
            \item Calcule la varianza de la variable aleatoria que modela el número de veces que hay que lanzar los dados hasta obtener 2 veces suma divisor de 2012.
        \end{enumerate}
    \end{ejercicio}

    \newpage
    \setcounter{ejercicio}{0} % Reiniciar contador de ejercicios
    \noindent
    % \textbf{Solución.}

\end{document}